\section{Robustness}
\label{sec:robust}

A parser that only works when everything goes right is not a parser. Three cases have to be
handled, and all three happen for real.

\subsection{Junk before SOF}

The node is connected in the middle of an ongoing transfer, or a disturbance adds bytes. The
stream looks like this:

\begin{console}
32 74 00 FF A5 F7 03 03 25 17 7F 05 10 20 30 02 C2
|---------| |-------------------------------------|
   junk                    the frame
\end{console}

The solution falls out of the state machine for free: in \code{WaitForSof1}, every byte that is
not \code{0xA5} is discarded. The parser stands still until the stream starts to look like a
frame.

One case is worth thinking through in particular: the parser is in \code{WaitForSof2} and gets a
byte that is not \code{0xF7}. The naive answer is to go back to \code{WaitForSof1} and discard the
byte. But if the byte was \code{0xA5}, you have just discarded the start of a real frame. The
right answer is to stay in \code{WaitForSof2}: the byte you just saw may be the first of the new
frame.

\subsection{Dropped bytes}

A byte disappears. Every following field is then shifted by one step, the length field reads a
data byte as the length, and the parser waits for the wrong number of bytes.

There is no elegant rescue, and there should not be one: the frame \emph{is} lost. The point is
that the parser should detect it and get back to a known state, not that it should guess. The
checksum fails, \code{extractFrame()} returns \code{false}, the caller calls \code{reset()}, and
the next \code{0xA5F7} in the stream starts over.

\subsection{An implausible length}

This is the case that is easy to miss, and the only one in the chapter that is a security bug
rather than mere error handling.

The length field is one byte and can therefore be up to 255. The parser's buffer holds
\code{MaxPayloadLen} bytes. If the parser receives a length greater than that, and then
obediently writes in that many data bytes, it writes past the end of its own buffer.

\begin{cppcode}
case State::WaitForLen:
{
    // Reject a length the buffer cannot hold, else the payload would overflow myBuf.
    if (MaxPayloadLen < byte) { return reportParseError(); }

    myBuf[myParsedBytes++] = byte;
    myState                = State::WaitForType;
    break;
}
\end{cppcode}

\begin{aside}
\note[Rule] Every length that comes from outside is to be checked against the size of the buffer
\emph{before} it is used, not after. The rule holds everywhere in this book, and it holds for
\code{TX_DLC} in \kapref{ch:driver} for exactly the same reason.
\end{aside}

\subsection{Two frames in a row}

When a frame is complete and extracted, the parser has to be reset before the next one can be
received. Forget that, and it stays in \code{Ready} and discards every new byte.

That is why \code{reset()} is in the public interface, and that is why the node in
\kapref{ch:bussar} calls it after every extracted frame.
